\documentclass[12pt]{article}
\usepackage[utf8]{inputenc}
\usepackage[T1]{fontenc}
\usepackage{hyperref}
\title{Folopa Language Resources for K-12 Curriculum Development}
\author{Compiled by Claude Sonnet 5.5 (claude-sonnet-5-5) for Andr\'as Kornai}
\date{}
\begin{document}
\maketitle

\section{Introduction}
Folopa (ISO 639-3 ppo; Glottocode folo1238) is a Papuan language of the small Teberan family, spoken in Gulf and Southern Highlands Provinces, Papua New Guinea. Speaker estimates in the results range from about 2,500 to 3,000, and one source lists a higher figure. Little material was found in the searches.

\section{Primary Educational Resources}
No primers, textbooks or school programmes in Folopa were found in the search results.

\section{Linguistic and Reference Materials}
\begin{itemize}
\item Wikipedia overview of Folopa: \url{https://en.wikipedia.org/wiki/Folopa_language}
\item Wikipedia on the Teberan family (Folopa and Dadibi): \url{https://en.wikipedia.org/wiki/Teberan_languages}
\item New Guinea World (University of Zurich) page on Folopa: \url{https://newguineaworld.linguistik.uzh.ch/families/papuan-gulf/tua-river/dadibi-folopa/folopa}
\end{itemize}

\section{Religious and Community Resources}
\begin{itemize}
\item Joshua Project entry under the name Podopa (a related name in the results; the match to Folopa was not verified): \url{https://joshuaproject.net/index.php/people_groups/14423}
\end{itemize}

\section{Digital and Audio Resources}
\begin{itemize}
\item Global Recordings Network, language page 1471 (listed under Folopa and its varieties; contents not inspected): \url{https://globalrecordings.net/en/language/1471}
\item Global Recordings Network, language page 9849 (contents not inspected): \url{https://globalrecordings.net/en/language/9849}
\end{itemize}

\section{Curriculum Development Context}
The results show no existing Folopa school materials, dictionary or orthography. A curriculum would need new community-led development; variety names differ between sources and should be checked with speakers and linguistic references.

\end{document}